\section{Lista de comprobación y mejora continua}

Antes de considerar terminada una funcionalidad, completar esta lista. Marcar ``no aplica'' solo cuando se anote una explicación breve.

\subsection{Interfaz y accesibilidad}
\begin{itemize}[leftmargin=*]
  \item El texto tiene contraste suficiente con su fondo y el contraste se ha comprobado con una herramienta.
  \item La información no depende solo del color, el sonido o un icono sin etiqueta.
  \item Los textos son legibles, directos y consistentes; las mayúsculas se limitan a rótulos breves.
  \item No hay ventanas emergentes automáticas o innecesarias; los avisos aparecen en contexto.
  \item Los errores explican cómo continuar y no borran datos introducidos sin avisar.
  \item Los controles funcionan con teclado y muestran un foco visible; se ha comprobado la ampliación del texto.
\end{itemize}

\subsection{Backend y calidad}
\begin{itemize}[leftmargin=*]
  \item El servidor valida los datos recibidos y aplica las reglas de negocio.
  \item Cada operación comprueba los permisos necesarios y devuelve solo los datos requeridos.
  \item Los errores se gestionan de forma controlada, sin exponer información interna ni marcar como guardado lo que falló.
  \item Se han probado el flujo normal y los casos de error relevantes.
  \item Otra persona ha revisado el cambio y los criterios de aceptación se pueden verificar.
  \item Las incidencias, decisiones y limitaciones pendientes están registradas y tienen seguimiento.
\end{itemize}

\subsection{Registro de formación y decisiones}
El equipo anotará la fecha de cada sesión, asistentes, ejercicios realizados, dudas que requieran seguimiento y cambios acordados para el manual. Para cada decisión relevante se registrará qué se decidió, por qué, quién la validó y si queda pendiente confirmación del centro. No se incluirán en actas ni ejemplos datos personales reales del alumnado.

\subsection{Referencias para ampliar la formación}
\begin{itemize}[leftmargin=*]
  \item W3C, \textit{Web Content Accessibility Guidelines (WCAG) 2.2}: \url{https://www.w3.org/TR/WCAG22/}.
  \item W3C WAI, comprensión del criterio de contraste mínimo: \url{https://www.w3.org/WAI/WCAG22/Understanding/contrast-minimum.html}.
  \item W3C WAI-ARIA Authoring Practices Guide, patrón de diálogo modal: \url{https://www.w3.org/WAI/ARIA/apg/patterns/dialog-modal/}.
\end{itemize}

Estas fuentes ayudan a concretar las comprobaciones técnicas; el equipo debe adaptarlas a las funciones reales del proyecto y contrastar la experiencia con sus usuarios previstos.